\begin{theorem}[Short tuples with prescribed arbitrary finite-group values]
\label{thm:group-typed-shortest}
Let $G=(V,\Sigma,P,S,\fan)$ be a reduced good MCFG
of fan-out at most $f$ and maximum rule rank
$r_{\max}\ge1$.  Let $g:\Sigma^*\to A$ be a
homomorphism into \emph{any} finite group $A$.
For every $B\in V$ and every realizable
$\mathbf s\in A^{\fan(B)}$ there is
$\tuple u\in L_B(G)$ with
$g^{(\fan(B))}(\tuple u)=\mathbf s$ and
\[
 \|\tuple u\|_1\le
 \beta_g^{\mathrm{grp}}:=
 (r_{\max}+1)\bar\theta_G
 \max\{1,(M_{A,f}-1)|V||A|^f\}.
\]
In particular $\beta_g^{\mathrm{grp}}
=O_{A,f,r_{\max}}(\|G\|\bar\theta_G)$.
\end{theorem}
\begin{proof}
Choose a terminal derivation from $B$ having the
prescribed root vector, lexicographically minimizing
(total length, node count).  Label each node by
its nonterminal and componentwise group value.
An ancestor and a proper descendant cannot have
the same label: replacing the ancestor by the
descendant would preserve every ancestral group
vector (by the ordered evaluation of rule templates),
would not increase total length (by nondeletion),
and would decrease node count.  Hence each
root-to-leaf path has at most
$n_g=|V||A|^f$ nodes.

Call a node heavy when its subtree's total yield
length exceeds $\bar\theta_G$.  Heavy nodes are
ancestor-closed.  By
Lemma~\ref{lem:group-heavy-antichain},
their leaves number at most $M_{A,f}-1$;
all heavy nodes lie on paths to those leaves,
so their number is at most
$(M_{A,f}-1)n_g$.
If the root is light, its yield already has
length at most $\bar\theta_G$.
Otherwise, each heavy rule contributes at most
$\bar\theta_G$ explicit terminal symbols, and
each of its at most $r_{\max}$ light child
subtrees contributes at most $\bar\theta_G$.
Counting the heavy nodes and their maximal light
children gives the displayed upper bound.
\end{proof}

\begin{lemma}[Short canonical tuples for window--arbitrary-group typing]
\label{lem:window-group-short}
Under the hypotheses of Theorem~\ref{thm:group-typed-shortest},
fix $k,\ell\ge0$ and a finite typing $h$ satisfying
$h\preceq(h_\Sigma^{k,\ell},g)$.
Let $\widetilde G_0$ be the trimmed full $h$-refinement
of $G$, and let $n_{\mathrm{typ}}$ denote its number of
non-start nonterminals.  Every non-start
$X$ of $\widetilde G_0$ admits a canonical tuple
$\omega(X)$ of total length at most
\[
 f(k+\ell+1)n_{\mathrm{typ}}
 \bigl(\bar\theta_G+r_{\max}\beta_g^{\mathrm{grp}}\bigr).
\]
\end{lemma}
\begin{proof}
Choose a terminal derivation from $X$ minimizing the
number of nodes.  Along a root-to-leaf path, a repeated
typed nonterminal could be bypassed, preserving the root
$h$-type and decreasing the number of nodes; thus
each path has at most $n_{\mathrm{typ}}$ nodes.
As in Lemma~\ref{lem:boundary-short-canonical},
mark the root positions determining the
$h_\Sigma^{k,\ell}$-value: all positions in a short
component, the first $k$ and last $\ell$ positions
in a long component, and one position when
$k=\ell=0$.  The union of ancestor paths to the
at most $f(k+\ell+1)$ marks is a skeleton with at most
$f(k+\ell+1)n_{\mathrm{typ}}$ nodes.

For each child subtree omitted from this skeleton,
record its underlying nonterminal $B$ and the
componentwise group vector actually realized.
Theorem~\ref{thm:group-typed-shortest} supplies
a replacement derivation from $B$ of length at most
$\beta_g^{\mathrm{grp}}$ and the \emph{same group
value in each component}.  Since $G$ is good,
all replacement components are nonempty.
No protected position changes, and each root
component that was long remains long; hence the
window value is preserved.  The ordered evaluation
of rule templates preserves every root group value
because each replaced child component retains its
group value; no commutativity is used.
Consequently the product $(h_\Sigma^{k,\ell},g)$
and the coarser root $h$-type are unchanged.
The underlying $G$-derivation lifts to the
trimmed full refinement with root $X$ by
Proposition~\ref{prop:output-invariants}(ii).

Each skeleton rule contributes at most
$\bar\theta_G$ explicit terminals and at most
$r_{\max}$ omitted children, each contributing at most
$\beta_g^{\mathrm{grp}}$ terminals.  The displayed
bound follows.
\end{proof}

\begin{theorem}[Polynomial characteristic data for window--finite-group typing]
\label{thm:window-group-polydata}
Fix $f,r_{\max},k,\ell,\Sigma$, any finite group
$A$, a homomorphism $g:\Sigma^*\to A$, and a
fixed finite typing $h$ satisfying
$h\preceq(h_\Sigma^{k,\ell},g)$.
For every reduced good fan-out-$\le f$,
rank-$\le r_{\max}$ grammar $G$ of a target in
$\Cmcf{f}{h}$, the capped reconstruction operator
$\mathcal B^{\le r_{\max}}_{f,h}$ has a characteristic
sample such that
\[
 |\Wit(\widetilde G_0)|=O(\|G\|),\qquad
 \|\Wit(\widetilde G_0)\|_+
 =O_{f,r_{\max},k,\ell,h,g,\Sigma}
   (\|G\|^3\bar\theta_G).
\]
\end{theorem}
\begin{proof}
By Lemma~\ref{lem:window-group-short},
canonical tuples in the trimmed refinement have
length at most
\[
 b_{\mathrm{can}}
 \le f(k+\ell+1)n_{\mathrm{typ}}
 (\bar\theta_G+r_{\max}\beta_g^{\mathrm{grp}})
 =O(\|G\|^2\bar\theta_G),
\]
since $n_{\mathrm{typ}}\le|V||M|^f=O(\|G\|)$
and the typing and rank bounds are fixed.
Lemma~\ref{lem:untyped-contexts} bounds each
characteristic word by
$2\bar\theta_G+r_{\max}b_{\mathrm{can}}
=O(\|G\|^2\bar\theta_G)$, and
Proposition~\ref{prop:witness-cardinality}
bounds the number of words by $O(\|G\|)$.
The optional empty-word witness adds only a
constant to the encoded total.
\end{proof}

\begin{remark}
This strengthening does not assert an original-thickness
polynomial for arbitrary finite \emph{monoid} typings.
For fixed $f$ the constant
$M_{A,f}=U_f\max\{3,|A|^f\}^{2^{f^2}}$ is polynomial in
$|A|$.  It replaces the multicolor Ramsey bound
$U_f\,c_f^{\,c_f(N_{A,f}-1)+1}$ with $c_f=2^{f^2}$ used in
earlier versions, which is exponential in $|A|^f$.  The
dependence on $f$ remains doubly exponential.  The abelian
argument of Appendix~\ref{app:window-abelian} gives a much
smaller constant, linear in $f$, when $A$ is commutative
(Remark~\ref{rem:davenport-size}).
\end{remark}
